\documentclass[11pt, french, openany]{report}

% preamble.tex — BTS FED GCF — Cours complet 2 ans
\usepackage[T1]{fontenc}
\usepackage[utf8]{inputenc}
\usepackage[french]{babel}
\usepackage{geometry}
\geometry{a4paper, margin=2.5cm}

% Mathématiques et unités
\usepackage{amsmath, amssymb, mathtools}
\usepackage{siunitx}
\sisetup{locale=FR, output-decimal-marker={,}}

% Tableaux et figures
\usepackage{booktabs, tabularx, multirow}
\usepackage{graphicx, float}
\usepackage{tikz}
\usetikzlibrary{shapes.geometric, arrows.meta, positioning}

% Mise en forme
\usepackage{xcolor}
\definecolor{btsprimary}{RGB}{0, 82, 147}
\definecolor{btssecondary}{RGB}{220, 237, 255}
\definecolor{btsaccent}{RGB}{255, 140, 0}

\usepackage{tcolorbox}
\tcbuselibrary{breakable, skins}

% Boîtes pédagogiques
\newtcolorbox{definition}[1][]{
  colback=btssecondary, colframe=btsprimary,
  fonttitle=\bfseries, title=Définition,#1
}
\newtcolorbox{methode}[1][]{
  colback=orange!10, colframe=btsaccent,
  fonttitle=\bfseries, title=Méthode,#1
}
\newtcolorbox{attention}[1][]{
  colback=red!8, colframe=red!60!black,
  fonttitle=\bfseries, title=Attention,#1
}
\newtcolorbox{exemple}[1][]{
  colback=green!8, colframe=green!50!black,
  fonttitle=\bfseries, title=Exemple résolu,#1
}

% En-têtes et pieds de page
\usepackage{fancyhdr}
\pagestyle{fancy}
\fancyhf{}
\fancyhead[L]{\small\textcolor{btsprimary}{BTS FED — Option GCF}}
\fancyhead[R]{\small\textcolor{btsprimary}{\leftmark}}
\fancyfoot[C]{\thepage}
\renewcommand{\headrulewidth}{0.4pt}

% Liens
\usepackage{hyperref}
\hypersetup{
  colorlinks=true, linkcolor=btsprimary,
  urlcolor=btsprimary, citecolor=btsprimary,
  pdftitle={Cours BTS FED — Option GCF},
  pdfauthor={Généré par agent Claude Code}
}

% Numérotation des exercices
\newcounter{exercice}[chapter]
\newcommand{\exercice}[1]{%
  \stepcounter{exercice}%
  \subsection*{Exercice \theexercice\ — #1}%
}


\begin{document}

% ============================================
% PAGE DE TITRE — LIVRET SITUATIONS RÉELLES
% Couleur identitaire : FedGreen
% ============================================

\begin{titlepage}
    \begin{tikzpicture}[remember picture, overlay]
        % Bande latérale gauche en FedGreen
        \node[anchor=north west, inner sep=0pt] at (current page.north west) {
            \begin{tcolorbox}[
                enhanced, colback=FedGreen, colframe=FedGreen, arc=0pt, outer arc=0pt,
                width=1.5cm, height=\paperheight, boxrule=0pt,
                left=0pt, right=0pt, top=0pt, bottom=0pt
            ]
            \end{tcolorbox}
        };

        % Texte rotaté sur la bande latérale
        \node[rotate=90, anchor=west, white] at ([xshift=0.85cm, yshift=2cm]current page.south west) {
            \sffamily\Huge\bfseries SITUATIONS RÉELLES
        };

        % Contenu principal
        \node[anchor=north west, text width=14cm] at ([xshift=3cm, yshift=-4cm]current page.north west) {
            {\huge\sffamily\color{FedBlue} \textbf{BTS FLUIDES ÉNERGIES DOMOTIQUE}}\\[0.5cm]
            {\large\sffamily OPTION A : GÉNIE CLIMATIQUE ET FLUIDIQUE}\\[2.5cm]

            {\fontsize{40}{45}\sffamily\bfseries\selectfont LIVRET}\\[0.3cm]
            {\fontsize{40}{45}\sffamily\bfseries\selectfont SITUATIONS}\\[0.3cm]
            {\fontsize{40}{45}\sffamily\bfseries\color{FedGreen}\selectfont RÉELLES}\\[0.8cm]
            {\color{FedGreen}\rule{10cm}{3pt}}\\[1cm]

            {\large\sffamily\color{FedBlue}\bfseries Projets de groupe — 2 à 3 étudiants}\\[0.4cm]
            {\normalsize\sffamily\color{gray}
              3 projets sur l'année $\cdot$ 2 notes par projet $\cdot$ Soutenance orale}
        };

        % Encadré récapitulatif des 3 SR
        \node[anchor=north west, text width=13cm] at ([xshift=3cm, yshift=-13.5cm]current page.north west) {
            \begin{tcolorbox}[
                enhanced, colback=FedGreen!8, colframe=FedGreen,
                arc=3pt, boxrule=1pt
            ]
            \sffamily\small
            \begin{tabular}{cll}
              \textbf{SR-1} & Réhabilitation thermique — Maison individuelle & Semestre 1 \\[4pt]
              \textbf{SR-2} & Climatisation et traitement d'air — Bureaux tertiaires & Semestre 2 \\[4pt]
              \textbf{SR-3} & Conception CVC complète — Immeuble résidentiel R+4 & Semestre 3 \\
            \end{tabular}
            \end{tcolorbox}
        };

        % Informations au bas
        \node[anchor=south west] at ([xshift=3cm, yshift=2cm]current page.south west) {
            \sffamily
            \begin{tabular}{ll}
                \textbf{Établissement :} & Lycée Gaudier Brzeska \\
                \textbf{Enseignant :} & M. Dylan PIMONT \\
                \textbf{Module :} & Projets professionnels BTS FED GCF
            \end{tabular}
        };

        % Décoration inférieure en FedGreen
        \fill[FedGreen] (current page.south west)
            rectangle ([yshift=8pt] current page.south east);
    \end{tikzpicture}
\end{titlepage}

\newpage

% ============================================
% TABLE DES MATIÈRES
% ============================================

\tableofcontents
\newpage

% ============================================
% INTRODUCTION AU LIVRET
% ============================================

\chapter*{Introduction — Comment utiliser ce livret ?}
\addcontentsline{toc}{chapter}{Introduction}

Ce livret regroupe les \textbf{trois Situations Réelles} du BTS FED option GCF.
Chaque situation est un projet professionnel complet, ancré dans le réel du secteur
du génie climatique et fluidique en France.

\begin{tcolorbox}[
  enhanced, colback=FedGreen!8, colframe=FedGreen, arc=3pt,
  fonttitle=\bfseries\sffamily, coltitle=white, colbacktitle=FedGreen,
  title={Distinction des 3 fascicules du cours BTS FED GCF}
]
\begin{center}
\begin{tikzpicture}[>=Stealth, thick, scale=0.95]

  % Fascicule 1 — Cours (bleu)
  \draw[FedBlue, very thick, fill=FedBlue!10, rounded corners=5pt]
    (0,0) rectangle (3.5,5.0);
  \fill[FedBlue] (0,0) rectangle (0.5,5.0);
  \node[white, rotate=90, font=\tiny\bfseries] at (0.25,2.5) {COURS};
  \node[FedBlue, font=\small\bfseries] at (2.0,4.2) {Fascicule 1};
  \node[FedBlue, font=\footnotesize] at (2.0,3.5) {\textbf{Référentiel de cours}};
  \node[FedBlue, font=\tiny, align=center] at (2.0,2.5)
    {12 chapitres\\Théorie + méthodes\\Exemples résolus};
  \node[FedBlue!60, font=\tiny] at (2.0,1.3) {main.tex};
  \node[FedBlue, font=\Large\bfseries] at (2.0,0.5) {\color{FedBlue}$\blacksquare$};

  % Fascicule 2 — Exercices (orange)
  \draw[FedOrange, very thick, fill=FedOrange!10, rounded corners=5pt]
    (4.5,0) rectangle (8.0,5.0);
  \fill[FedOrange] (4.5,0) rectangle (5.0,5.0);
  \node[white, rotate=90, font=\tiny\bfseries] at (4.75,2.5) {EXOS};
  \node[FedOrange, font=\small\bfseries] at (6.25,4.2) {Fascicule 2};
  \node[FedOrange, font=\footnotesize] at (6.25,3.5) {\textbf{Livret d'exercices}};
  \node[FedOrange, font=\tiny, align=center] at (6.25,2.5)
    {12 séries d'exercices\\Corrigés détaillés\\Progression niveau A→C};
  \node[FedOrange!60, font=\tiny] at (6.25,1.3) {livret-exercices.tex};
  \node[FedOrange, font=\Large\bfseries] at (6.25,0.5) {\color{FedOrange}$\blacksquare$};

  % Fascicule 3 — Situations (vert)
  \draw[FedGreen, very thick, fill=FedGreen!10, rounded corners=5pt]
    (9.0,0) rectangle (12.5,5.0);
  \fill[FedGreen] (9.0,0) rectangle (9.5,5.0);
  \node[white, rotate=90, font=\tiny\bfseries] at (9.25,2.5) {SR};
  \node[FedGreen, font=\small\bfseries] at (10.75,4.2) {Fascicule 3};
  \node[FedGreen, font=\footnotesize] at (10.75,3.5) {\textbf{Situations Réelles}};
  \node[FedGreen, font=\tiny, align=center] at (10.75,2.5)
    {3 projets annuels\\Groupe 2-3 étudiants\\2 notes + soutenance};
  \node[FedGreen!60, font=\tiny] at (10.75,1.3) {livret-situations.tex};
  \node[FedGreen, font=\Large\bfseries] at (10.75,0.5) {\color{FedGreen}$\blacksquare$};

  % Flèches progression
  \draw[->, gray, thick] (3.6,2.5) -- (4.4,2.5)
    node[midway, above, font=\tiny, gray] {alimente};
  \draw[->, gray, thick] (8.1,2.5) -- (8.9,2.5)
    node[midway, above, font=\tiny, gray] {mobilise};

\end{tikzpicture}
\end{center}
\end{tcolorbox}

\subsection*{Modalités d'évaluation des Situations Réelles}

\begin{center}
\begin{tabular}{clccc}
  \toprule
  \textbf{SR} & \textbf{Titre} & \textbf{Semestre} & \textbf{Note 1} & \textbf{Note 2} \\
  \midrule
  SR-1 & Réhabilitation maison individuelle & 1 & Dossier (coeff.\ 3) & Oral (coeff.\ 2) \\
  SR-2 & Tertiaire — CTA + GTB & 2 & Dossier (coeff.\ 3) & Oral + démo (coeff.\ 2) \\
  SR-3 & Immeuble R+4 — Conception complète & 3 & Dossier (coeff.\ 3) & Commission (coeff.\ 2) \\
  \bottomrule
\end{tabular}
\end{center}

\begin{attention}[Travail en groupe — Règles]
  \begin{itemize}
    \item Groupes de \textbf{2 à 3 étudiants maximum} — constitués en semaine 1
    \item \textbf{Répartition du travail} documentée dans la première page du dossier
    \item Les questions individuelles lors de la soutenance évaluent chaque étudiant séparément
    \item Un plagiat ou dossier identique entre groupes entraîne la note de \textbf{0} pour les deux groupes
  \end{itemize}
\end{attention}

\newpage

% ============================================
% SITUATIONS RÉELLES
% ============================================

% ================================================================
%  Situation Réelle 1 — Réhabilitation thermique d'une maison individuelle
%  BTS FED — Option GCF — Semestre 1
%  Chapitres mobilisés : 01 (thermique), 02 (hydraulique), 04 (production), 05 (émission)
% ================================================================

\chapter{SR-1 — Réhabilitation thermique d'une maison individuelle}

\begin{tcolorbox}[
  enhanced, colback=FedGreen!8, colframe=FedGreen, arc=3pt,
  fonttitle=\bfseries\sffamily\large, coltitle=white, colbacktitle=FedGreen,
  title={Situation Réelle 1 — Fiche d'identité du projet},
  width=\textwidth
]
\begin{tabular}{@{}ll@{\hspace{2cm}}ll@{}}
  \textbf{Semestre} & 1 &
  \textbf{Durée} & 6 semaines \\[4pt]
  \textbf{Groupe} & 2 à 3 étudiants &
  \textbf{Chapitres} & 01, 02, 04, 05 \\[4pt]
  \textbf{Note 1} & Dossier technique (coeff.\ 3) &
  \textbf{Note 2} & Soutenance orale (coeff.\ 2) \\[4pt]
  \textbf{Compétences} & C2-3, C3-3, C5-1, C8-3, C8-4 & & \\
\end{tabular}
\end{tcolorbox}

\bigskip

% ================================================================
\section{Contexte professionnel}
% ================================================================

Vous êtes technicien supérieur GCF dans l'entreprise \textbf{ThermoPro SARL} (Grenoble, 38).
Un particulier, M.\ et Mme Aubert, souhaite rénover leur maison individuelle construite
en \textbf{1975} afin d'améliorer son étiquette DPE de \textbf{D} vers \textbf{B} et
de remplacer la chaudière fioul existante par une \textbf{pompe à chaleur air/eau}.

\subsection*{Description du bâtiment}

\begin{figure}[htbp]
  \centering
  \begin{tikzpicture}[scale=1.0, >=Stealth]

    % --- Sol ---
    \fill[gray!20] (-0.3,-0.3) rectangle (10.3, 0);
    \draw[thick] (-0.3,0) -- (10.3,0);

    % --- Murs extérieurs ---
    \draw[thick, FedBlue, fill=FedBlue!10]
      (0,0) rectangle (10,4.5);

    % --- Toit ---
    \draw[thick, FedBlue!80, fill=FedBlue!20]
      (-0.5,4.5) -- (5,7) -- (10.5,4.5) -- cycle;

    % --- Plancher intermédiaire (R+1) ---
    \draw[FedOrange, thick] (0,2.2) -- (10,2.2);
    \node[FedOrange, font=\footnotesize] at (5,2.0) {Plancher intermédiaire};

    % --- Fenêtres RDC ---
    \draw[FedBlue!60, fill=cyan!15, thick] (1,0.5) rectangle (2.5,1.8);
    \draw[FedBlue!60, fill=cyan!15, thick] (4,0.5) rectangle (5.8,1.8);
    \draw[FedBlue!60, fill=cyan!15, thick] (7.2,0.5) rectangle (8.8,1.8);
    % --- Porte ---
    \draw[brown!70, fill=brown!20, thick] (3.1,0) rectangle (3.9,2.0);
    \draw[brown!70] (3.1,1.0) arc (180:0:0.4);

    % --- Fenêtres R+1 ---
    \draw[FedBlue!60, fill=cyan!15, thick] (1.5,2.6) rectangle (3,3.8);
    \draw[FedBlue!60, fill=cyan!15, thick] (5,2.6) rectangle (6.5,3.8);
    \draw[FedBlue!60, fill=cyan!15, thick] (7.5,2.6) rectangle (9,3.8);

    % --- Velux ---
    \draw[FedBlue!60, fill=cyan!20, thick] (3.5,5.2) rectangle (4.8,6.0);
    \draw[FedBlue!60, fill=cyan!20, thick] (5.5,5.2) rectangle (6.8,6.0);

    % --- Annotations ---
    \node[font=\small\bfseries, FedBlue] at (5, 1.0) {RDC — Surface : \SI{60}{\metre\squared}};
    \node[font=\small\bfseries, FedBlue] at (5, 3.3) {R+1 — Surface : \SI{60}{\metre\squared}};
    \node[font=\small, FedBlue!80] at (5, 5.7) {Combles perdus};

    % --- Côtes ---
    \draw[<->, gray] (-0.1, 0) -- (-0.1, 4.5)
      node[midway, left, font=\footnotesize] {\SI{5.5}{\metre}};
    \draw[<->, gray] (0, -0.2) -- (10, -0.2)
      node[midway, below, font=\footnotesize] {\SI{10}{\metre}};

    % --- Isolation existante (hachures) ---
    \foreach \y in {0.3, 0.7, 1.1, 1.5, 1.9, 2.4, 2.8, 3.2, 3.6, 4.0, 4.3}{
      \draw[red!40, thin] (0,\y) -- (0.25, \y+0.2);
      \draw[red!40, thin] (9.75,\y) -- (10,\y+0.2);
    }
    \node[red!60, font=\tiny, rotate=90] at (0.15, 2.5) {Laine verre 60mm};

    % --- Chaudière fioul existante ---
    \draw[FedRed, fill=FedRed!15, thick, rounded corners=2pt]
      (0.5, -2.5) rectangle (2.5, -1.0);
    \node[font=\footnotesize\bfseries, FedRed] at (1.5,-1.75) {Chaudière fioul};
    \node[font=\tiny, FedRed] at (1.5,-2.1) {18 kW — η = 82\%};
    \draw[FedRed, thick, ->] (1.5,-1.0) -- (1.5,-0.0);

    % --- PAC future ---
    \draw[FedGreen, fill=FedGreen!15, thick, rounded corners=2pt, dashed]
      (7, -2.5) rectangle (9.5, -1.0);
    \node[font=\footnotesize\bfseries, FedGreen] at (8.25,-1.6) {PAC air/eau};
    \node[font=\tiny, FedGreen] at (8.25,-2.1) {à dimensionner};
    \node[font=\tiny, FedGreen] at (8.25,-2.35) {A7/W35};

    % --- Légende ---
    \draw[FedRed, thick] (0.5,-3.2) -- (1.0,-3.2);
    \node[right, font=\tiny] at (1.0,-3.2) {Existant (dépose)};
    \draw[FedGreen, thick, dashed] (3.5,-3.2) -- (4.0,-3.2);
    \node[right, font=\tiny] at (4.0,-3.2) {Projet};
    \draw[red!40] (6.5,-3.2) -- (7.0,-3.2);
    \node[right, font=\tiny] at (7.0,-3.2) {Isolation existante};

  \end{tikzpicture}
  \caption{Coupe schématique de la maison Aubert — Grenoble (H1c) — Construction 1975
    — S\textsubscript{habitable} = \SI{120}{\metre\squared} — Chaudière fioul à remplacer par PAC air/eau}
  \label{fig:sr1-coupe-maison}
\end{figure}

\subsection*{Données du bâtiment existant}

\begin{center}
\begin{tabular}{lllr}
  \toprule
  \textbf{Élément} & \textbf{Composition actuelle} & \textbf{U actuel} & \textbf{Proportion} \\
  \midrule
  Murs extérieurs & Parpaings 20 cm + enduits & \SI{1.8}{\watt\per\metre\squared\per\kelvin} & \SI{140}{\metre\squared} \\
  Plancher bas & Dalle béton sur terre-plein & \SI{0.6}{\watt\per\metre\squared\per\kelvin} & \SI{60}{\metre\squared} \\
  Toiture & Tuiles + laine de verre 60 mm & \SI{0.55}{\watt\per\metre\squared\per\kelvin} & \SI{65}{\metre\squared} \\
  Fenêtres RDC & Simple vitrage bois & \SI{4.5}{\watt\per\metre\squared\per\kelvin} & \SI{9}{\metre\squared} \\
  Fenêtres R+1 & Simple vitrage bois & \SI{4.5}{\watt\per\metre\squared\per\kelvin} & \SI{9}{\metre\squared} \\
  Porte d'entrée & Bois plein & \SI{3.2}{\watt\per\metre\squared\per\kelvin} & \SI{2}{\metre\squared} \\
  \midrule
  \textbf{Pont thermique linéique} & Jonction mur/plancher & $\psi = \SI{0.8}{\watt\per\metre\per\kelvin}$ & \SI{40}{\metre} \\
  \bottomrule
\end{tabular}
\end{center}

\begin{definition}[Données climatiques — Grenoble, zone H1c]
  Température extérieure de base : $T_{ext,base} = \SI{-11}{\celsius}$ (NF EN 12831) \\
  Température intérieure de confort : $T_{int} = \SI{19}{\celsius}$ \\
  Degré-jours unifiés (DJU) : \textbf{2 900 DJU/an} (référence 18°C) \\
  Ensoleillement annuel : \SI{1950}{\hour\per\year}
\end{definition}

% ================================================================
\section{Schéma de principe de l'installation projetée}
% ================================================================

\begin{figure}[htbp]
  \centering
  \begin{tikzpicture}[
    >=Stealth, thick,
    bloc/.style={draw, rounded corners=3pt, minimum width=2.8cm, minimum height=1.0cm,
                 align=center, font=\small},
    tuyau/.style={draw, ->, thick},
    dep/.style={FedOrange, thick},   % départ chaud
    ret/.style={FedBlue, thick},     % retour froid
  ]

    % === PAC air/eau ===
    \node[bloc, fill=FedGreen!15, draw=FedGreen] (pac) at (0,0)
      {PAC air/eau\\{\footnotesize A7/W35}};

    % === Ballon tampon ===
    \node[bloc, fill=FedBlue!10, draw=FedBlue] (tampon) at (4,0)
      {Ballon\\tampon\\{\footnotesize 300 L}};

    % === Collecteur départ ===
    \node[bloc, fill=FedOrange!15, draw=FedOrange, minimum width=3.5cm] (coll) at (8,0)
      {Collecteur\\départ/retour};

    % === Émetteurs ===
    \node[bloc, fill=gray!10, draw=gray] (pch-rdc) at (11, 2.5)
      {PCH RDC\\{\footnotesize 60 m²}};
    \node[bloc, fill=gray!10, draw=gray] (rad-r1) at (11, -2.5)
      {Radiateurs R+1\\{\footnotesize BT 45/40°C}};

    % === Vase d'expansion ===
    \node[draw, circle, minimum size=0.8cm, fill=yellow!20, font=\tiny] (vase) at (4, 2.5)
      {VE};

    % === Pompes de circulation ===
    \node[draw, circle, minimum size=0.6cm, fill=gray!20, font=\tiny] (pompe1) at (2, 0.4)
      {P1};
    \node[draw, circle, minimum size=0.6cm, fill=gray!20, font=\tiny] (pompe2) at (6, 0.4)
      {P2};

    % === Connexions départ (orange) ===
    \draw[dep, ->] (pac.east) -- node[above, font=\tiny] {55°C} (tampon.west);
    \draw[dep, ->] (tampon.east) -- node[above, font=\tiny] {45°C} (coll.west);
    \draw[dep, ->] (coll.north east) |- (pch-rdc.west);
    \draw[dep, ->] (coll.south east) |- (rad-r1.west);

    % === Connexions retour (bleu) ===
    \draw[ret, ->] (pch-rdc.south) -- ++(0,-0.8) -| (coll.north);
    \draw[ret, ->] (rad-r1.north) -- ++(0,0.8) -| (coll.south);
    \draw[ret, ->] (coll.west|-tampon) -- node[below, font=\tiny] {35°C} (tampon.east|-tampon);
    \draw[ret, ->] (tampon.west) -- node[below, font=\tiny] {40°C} (pac.east|-pac.south);

    % === Vase d'expansion connexion ===
    \draw[gray, dashed] (vase) -- (tampon.north);

    % === Légende ===
    \draw[dep, ->] (0,-4.2) -- (1.2,-4.2) node[right, font=\tiny, black] {Circuit départ chaud};
    \draw[ret, ->] (3.5,-4.2) -- (4.7,-4.2) node[right, font=\tiny, black] {Circuit retour froid};
    \node[draw, circle, minimum size=0.5cm, fill=gray!20, font=\tiny] at (7.5,-4.2) {P};
    \node[right, font=\tiny] at (7.8,-4.2) {Circulateur};
    \node[draw, circle, minimum size=0.5cm, fill=yellow!20, font=\tiny] at (9.5,-4.2) {VE};
    \node[right, font=\tiny] at (9.8,-4.2) {Vase expansion};

    % === Sonde extérieure ===
    \node[draw, diamond, fill=purple!15, minimum size=0.5cm, font=\tiny] (sonde) at (-2, 2.0)
      {$T_e$};
    \draw[purple, dashed, ->] (sonde) -- (pac.north west);
    \node[purple, font=\tiny, above] at (-1, 2.0) {Régulation};
    \node[purple, font=\tiny, above] at (-1, 1.6) {climatique};

  \end{tikzpicture}
  \caption{Schéma hydraulique de principe — Installation PAC air/eau + ballon tampon + PCH RDC + radiateurs R+1
    — Régulation climatique par sonde extérieure}
  \label{fig:sr1-schema-hydraulique}
\end{figure}

% ================================================================
\section{Cahier des charges}
% ================================================================

\begin{methode}[Livrables attendus — Note 1 : Dossier technique écrit]
\textbf{À remettre en semaine 6 — Format PDF, 25 pages minimum}

\begin{enumerate}
  \item \textbf{Bilan thermique} selon NF EN 12831 :
    \begin{itemize}
      \item Calcul du coefficient $U$ de chaque paroi (formule de la résistance thermique équivalente)
      \item Calcul des déperditions totales $\Phi_{dep}$ du bâtiment [\si{\watt}]
      \item Calcul de $U_{bat}$ et comparaison avec le seuil RE2020
      \item Identification des travaux d'isolation complémentaires recommandés
    \end{itemize}

  \item \textbf{Dimensionnement PAC air/eau} :
    \begin{itemize}
      \item Calcul de la puissance utile nécessaire en régime de base et en point de bivalence
      \item Sélection d'un modèle constructeur (Daikin, Atlantic, Bosch) avec justification
      \item Calcul du COP en condition A7/W35 et du SCOP annuel (méthode bin)
      \item Dimensionnement du ballon tampon (volume, pertes thermiques)
    \end{itemize}

  \item \textbf{Réseau hydraulique} :
    \begin{itemize}
      \item Calcul des débits par circuit (plancher chauffant + radiateurs)
      \item Calcul des pertes de charge (tuyauteries + émetteurs + vannes)
      \item Sélection des circulateurs (courbe pompe/réseau, classe ErP)
      \item Dimensionnement du vase d'expansion (DTU 65.11)
    \end{itemize}

  \item \textbf{Émetteurs} :
    \begin{itemize}
      \item Plancher chauffant RDC : calcul densité de flux, pas de pose, pertes de charge boucles
      \item Radiateurs R+1 : sélection en régime basse température (45/40°C)
    \end{itemize}

  \item \textbf{Aspects économiques} :
    \begin{itemize}
      \item Devis estimatif par poste (€ HT)
      \item Calcul des économies annuelles (fioul → PAC, sur base DJU)
      \item Calcul du TRI avec et sans aides (MaPrimeRénov', éco-PTZ)
    \end{itemize}
\end{enumerate}
\end{methode}

\begin{methode}[Livrables attendus — Note 2 : Soutenance orale]
\textbf{Durée : 20 min exposé + 10 min questions — Semaine 7}

\begin{itemize}
  \item Support visuel : diaporama (10 à 15 slides) + schéma hydraulique A3 imprimé
  \item Chaque membre du groupe présente une partie du dossier
  \item Questions individuelles obligatoires posées par le jury sur les calculs
  \item Critères d'évaluation : clarté technique, rigueur des calculs, justification des choix
\end{itemize}
\end{methode}

% ================================================================
\section{Grille d'évaluation}
% ================================================================

\begin{center}
\begin{tabular}{p{5cm}ccp{5cm}}
  \toprule
  \textbf{Critère} & \textbf{Barème} & \textbf{Note} & \textbf{Indicateurs de réussite} \\
  \midrule
  \multicolumn{4}{l}{\textbf{Note 1 — Dossier technique (20 pts)}} \\
  \midrule
  Bilan thermique & 5 pts & /5 & U corrects, déperditions cohérentes ±10\% \\
  Dim. PAC & 4 pts & /4 & Puissance justifiée, modèle sélectionné \\
  Réseau hydraulique & 4 pts & /4 & Débits, $\Delta P$, circulateur justifiés \\
  Émetteurs & 3 pts & /3 & PCH et radiateurs dimensionnés \\
  Aspect économique & 2 pts & /2 & TRI calculé, aides identifiées \\
  Qualité rédactionnelle & 2 pts & /2 & Plan, syntaxe, figures légendées \\
  \midrule
  \multicolumn{4}{l}{\textbf{Note 2 — Soutenance orale (20 pts)}} \\
  \midrule
  Clarté de l'exposé & 5 pts & /5 & Structure logique, gestion du temps \\
  Maîtrise technique & 8 pts & /8 & Réponses aux questions individuelles \\
  Support visuel & 4 pts & /4 & Schémas clairs, lisibles \\
  Travail d'équipe & 3 pts & /3 & Répartition équilibrée \\
  \bottomrule
\end{tabular}
\end{center}

% ================================================================
\section{Ressources documentaires}
% ================================================================

\begin{itemize}
  \item \textbf{Normes :} NF EN 12831 (calcul besoins chaleur), NF EN 14511 (PAC), DTU 65.11 (sécurités)
  \item \textbf{Réglementation :} RE2020 — arrêté du 4 août 2021, Règlement (UE) 813/2013 (PAC)
  \item \textbf{Aides financières :} MaPrimeRénov' (ANAH), éco-PTZ, TVA 5,5\%
  \item \textbf{Logiciels conseillés :} Pleiades+Comfie (bilan thermique), ClimaWin (PAC)
  \item \textbf{Catalogues constructeurs :} Daikin Altherma 3, Atlantic Extensa+, Buderus Logatherm
  \item \textbf{Cours :} Chapitres 01 (thermique), 02 (hydraulique), 04 (production), 05 (distribution)
\end{itemize}

\finchapitre


% ================================================================
%  Situation Réelle 2 — Climatisation et traitement d'air d'un bâtiment tertiaire
%  BTS FED — Option GCF — Semestre 2
%  Chapitres mobilisés : 03 (aéraulique), 06 (climatisation), 07 (GTB), 08 (efficacité)
% ================================================================

\chapter{SR-2 — Climatisation et traitement d'air d'un bâtiment tertiaire}

\begin{tcolorbox}[
  enhanced, colback=FedGreen!8, colframe=FedGreen, arc=3pt,
  fonttitle=\bfseries\sffamily\large, coltitle=white, colbacktitle=FedGreen,
  title={Situation Réelle 2 — Fiche d'identité du projet},
  width=\textwidth
]
\begin{tabular}{@{}ll@{\hspace{2cm}}ll@{}}
  \textbf{Semestre} & 2 &
  \textbf{Durée} & 6 semaines \\[4pt]
  \textbf{Groupe} & 2 à 3 étudiants &
  \textbf{Chapitres} & 03, 06, 07, 08 \\[4pt]
  \textbf{Note 1} & Dossier technique (coeff.\ 3) &
  \textbf{Note 2} & Soutenance + démo GTB (coeff.\ 2) \\[4pt]
  \textbf{Compétences} & C2-1, C3-3, C5-2, C6-3, C7-4, C8-4 & & \\
\end{tabular}
\end{tcolorbox}

\bigskip

% ================================================================
\section{Contexte professionnel}
% ================================================================

Vous êtes bureau d'études chez \textbf{AirTech Ingénierie} (Lyon, 69).
Un promoteur immobilier vous confie la conception des installations CVC d'un
\textbf{immeuble de bureaux} neuf situé à \textbf{Lyon Part-Dieu} (zone H1b),
soumis à la réglementation \textbf{RE2020}.
Le bâtiment est classé \textbf{ERP de catégorie 3}.

\subsection*{Description du bâtiment}

\begin{figure}[htbp]
  \centering
  \begin{tikzpicture}[scale=0.85, >=Stealth, thick]

    % ============= PLAN RDC =============
    \node[font=\bfseries\small, FedBlue] at (5, 9.8) {Plan du plateau type — Niveaux R+1 et R+2};

    % Contour extérieur
    \draw[FedBlue, very thick, fill=FedBlue!5]
      (0,0) rectangle (10,8);

    % Couloir central
    \draw[FedOrange, thick, fill=FedOrange!10]
      (0,3.5) rectangle (10,4.5);
    \node[FedOrange, font=\footnotesize] at (5,4.0) {Couloir de circulation};

    % Bureau open-space gauche
    \draw[FedBlue!50, fill=FedBlue!8]
      (0.2,4.7) rectangle (4.8,7.8);
    \node[FedBlue, font=\footnotesize\bfseries] at (2.5,6.5) {Open Space A};
    \node[FedBlue, font=\tiny] at (2.5,6.0) {\SI{216}{\metre\squared}};
    \node[FedBlue, font=\tiny] at (2.5,5.6) {40 postes};

    % Bureau open-space droit
    \draw[FedBlue!50, fill=FedBlue!8]
      (5.2,4.7) rectangle (9.8,7.8);
    \node[FedBlue, font=\footnotesize\bfseries] at (7.5,6.5) {Open Space B};
    \node[FedBlue, font=\tiny] at (7.5,6.0) {\SI{216}{\metre\squared}};
    \node[FedBlue, font=\tiny] at (7.5,5.6) {40 postes};

    % Salles de réunion bas
    \draw[FedGreen!60, fill=FedGreen!10]
      (0.2,0.2) rectangle (2.8,3.3);
    \node[FedGreen, font=\tiny\bfseries] at (1.5,1.75) {Salle Réunion 1};
    \node[FedGreen, font=\tiny] at (1.5,1.3) {\SI{35}{\metre\squared}};

    \draw[FedGreen!60, fill=FedGreen!10]
      (3.0,0.2) rectangle (5.5,3.3);
    \node[FedGreen, font=\tiny\bfseries] at (4.25,1.75) {Salle Réunion 2};
    \node[FedGreen, font=\tiny] at (4.25,1.3) {\SI{38}{\metre\squared}};

    \draw[FedGreen!60, fill=FedGreen!10]
      (5.7,0.2) rectangle (7.8,3.3);
    \node[FedGreen, font=\tiny\bfseries] at (6.75,1.75) {Salle Réunion 3};
    \node[FedGreen, font=\tiny] at (6.75,1.3) {\SI{31}{\metre\squared}};

    % Local technique
    \draw[FedRed!60, fill=FedRed!10]
      (8.0,0.2) rectangle (9.8,3.3);
    \node[FedRed, font=\tiny\bfseries] at (8.9,1.75) {Local\\technique};
    \node[FedRed, font=\tiny] at (8.9,1.2) {CTA + GTB};

    % Bouches de soufflage (schématiques)
    \foreach \x in {1.0, 2.0, 3.0, 3.8, 6.0, 7.0, 8.0, 9.0}{
      \draw[cyan!60, fill=cyan!30] (\x, 7.5) circle (0.12);
    }
    \foreach \x in {1.0, 2.0, 3.0, 3.8, 6.0, 7.0, 8.0, 9.0}{
      \draw[orange!60, fill=orange!30] (\x, 5.2) circle (0.12);
    }

    % Légende bouches
    \draw[cyan!60, fill=cyan!30] (0.3, -0.6) circle (0.12);
    \node[right, font=\tiny] at (0.45,-0.6) {Bouche soufflage};
    \draw[orange!60, fill=orange!30] (3.0, -0.6) circle (0.12);
    \node[right, font=\tiny] at (3.15,-0.6) {Bouche reprise};

    % Côtes
    \draw[<->, gray] (-0.2, 0) -- (-0.2, 8)
      node[midway, left, font=\tiny] {\SI{20}{\metre}};
    \draw[<->, gray] (0, -0.2) -- (10, -0.2)
      node[midway, below, font=\tiny] {\SI{40}{\metre}};

  \end{tikzpicture}
  \caption{Plan du plateau de bureaux type (R+1 et R+2) — Surface totale par niveau : \SI{800}{\metre\squared}
    — Local technique en angle pour la CTA — Lyon Part-Dieu, zone H1b}
  \label{fig:sr2-plan-plateau}
\end{figure}

\subsection*{Données du bâtiment}

\begin{center}
\begin{tabular}{lrl}
  \toprule
  \textbf{Caractéristique} & \textbf{Valeur} & \textbf{Unité / Remarque} \\
  \midrule
  Surface totale chauffée/climatisée & \num{1600} & \si{\metre\squared} (2 niveaux × 800) \\
  Hauteur sous plafond & \num{2.8} & \si{\metre} \\
  Occupation maximale & 80 + 40 & personnes (open-space + salles) \\
  Charge thermique interne & \num{40} & \si{\watt\per\metre\squared} (bureaux RE2020) \\
  Vitrage (façades E et O) & \num{35} & \% de la surface de façade \\
  Facteur solaire $g$ (vitrage) & \num{0.28} & — (traitement anti-solaire) \\
  Renouvellement d'air neuf & \num{25} & \si{\metre\cubed\per\hour\per\person} (NF EN 16798) \\
  Température intérieure été & \num{26} & \si{\celsius} \\
  Température intérieure hiver & \num{19} & \si{\celsius} \\
  \bottomrule
\end{tabular}
\end{center}

\begin{attention}[Exigence maître d'ouvrage]
  Le maître d'ouvrage impose \textbf{obligatoirement} :
  \begin{itemize}
    \item Free-cooling aéraulique économiseur (bypass CTA si $T_{ext} < T_{int} - 2°C$)
    \item Récupération de chaleur sur air extrait : rendement $\eta \geq 85\,\%$ (échangeur rotatif)
    \item GTB de niveau 4 (BACS class C, EN 15232) avec supervision Web
    \item Cep\textsubscript{nr} $\leq \SI{50}{\kilo\watt\hour_{ep}\per\metre\squared\per\year}$ (RE2020 tertiaire)
  \end{itemize}
\end{attention}

% ================================================================
\section{Schéma de principe CTA double flux}
% ================================================================

\begin{figure}[htbp]
  \centering
  \begin{tikzpicture}[
    >=Stealth, thick,
    composant/.style={draw, minimum width=2.0cm, minimum height=1.2cm,
                      align=center, font=\scriptsize, rounded corners=2pt},
    flux_s/.style={FedBlue, thick, ->},   % air soufflage
    flux_e/.style={FedOrange, thick, ->}, % air extraction
  ]

    % ============ LIGNE SOUFFLAGE (haut) ============
    % Prise d'air neuf
    \node[composant, fill=cyan!10, draw=cyan!60] (pan) at (0, 3.5)
      {Prise air\\neuf\\{\tiny filtre G4}};

    % Échangeur récupérateur
    \node[composant, fill=purple!10, draw=purple!60, minimum width=2.5cm, minimum height=2.5cm]
      (echangeur) at (3.5, 2.5) {Échangeur\\rotatif\\{\tiny $\eta \geq 85\,\%$}};

    % Filtre fin
    \node[composant, fill=gray!10] (filtre) at (6.5, 3.5)
      {Filtre\\F7};

    % Batterie froide
    \node[composant, fill=FedBlue!20, draw=FedBlue] (batt-f) at (9, 3.5)
      {Batterie\\froide\\{\tiny groupe froid}};

    % Batterie chaude
    \node[composant, fill=FedOrange!20, draw=FedOrange] (batt-c) at (11.5, 3.5)
      {Batterie\\chaude\\{\tiny PAC / chaudière}};

    % Ventilateur soufflage
    \node[draw, circle, minimum size=0.9cm, fill=FedBlue!15, font=\tiny] (vs) at (13.8, 3.5)
      {$V_S$};

    % Réseau soufflage
    \node[composant, fill=FedBlue!8, draw=FedBlue!40, minimum width=2.2cm] (res-s) at (15.8, 3.5)
      {Réseau\\soufflage\\{\tiny bouches}};

    % ============ LIGNE EXTRACTION (bas) ============
    % Réseau reprise
    \node[composant, fill=FedOrange!8, draw=FedOrange!40, minimum width=2.2cm] (res-r) at (15.8, 1.5)
      {Réseau\\reprise\\{\tiny bouches}};

    % Ventilateur extraction
    \node[draw, circle, minimum size=0.9cm, fill=FedOrange!15, font=\tiny] (ve) at (13.8, 1.5)
      {$V_E$};

    % Bypass free-cooling
    \node[composant, fill=FedGreen!10, draw=FedGreen, minimum width=1.5cm] (bypass) at (11.0, 1.5)
      {Bypass\\free-cooling\\{\tiny registre motorisé}};

    % Rejet air vicié
    \node[composant, fill=gray!10, draw=gray] (rejet) at (3.5, 0.5)
      {Rejet\\air vicié\\{\tiny filtre antipollen}};

    % ============ CONNEXIONS SOUFFLAGE (bleu) ============
    \draw[flux_s] (pan.east) -- (echangeur.north west);
    \draw[flux_s] (echangeur.north east) -- (filtre.west);
    \draw[flux_s] (filtre.east) -- (batt-f.west);
    \draw[flux_s] (batt-f.east) -- (batt-c.west);
    \draw[flux_s] (batt-c.east) -- (vs.west);
    \draw[flux_s] (vs.east) -- (res-s.west);

    % ============ CONNEXIONS EXTRACTION (orange) ============
    \draw[flux_e] (res-r.west) -- (ve.east);
    \draw[flux_e] (ve.west) -- (bypass.east);
    \draw[flux_e] (bypass.west) -- (echangeur.south east);
    \draw[flux_e] (echangeur.south west) -- (rejet.east);

    % ============ BYPASS FREE-COOLING (vert) ============
    \draw[FedGreen, thick, dashed, ->]
      (bypass.north) |- node[pos=0.3, above, font=\tiny, FedGreen]
      {$T_{ext} < T_{int}-2°C$} (batt-f.south);

    % ============ ANNOTATIONS TEMPÉRATURES ============
    \node[font=\tiny, cyan!70!black] at (1.5, 4.2) {$T_{ext} = -11$ à $+38°C$};
    \node[font=\tiny, FedBlue] at (7.5, 4.2) {$T_{souf} = 16°C$ (été)};
    \node[font=\tiny, FedBlue] at (7.5, 3.85) {$T_{souf} = 22°C$ (hiver)};
    \node[font=\tiny, FedOrange!70] at (1.5, 0.2) {$T_{rej} \approx T_{ext} + 5°C$};
    \node[font=\tiny, FedOrange] at (11.0, 0.9) {$T_{rep} = 26°C$ (été)};

    % ============ LÉGENDE ============
    \draw[flux_s] (0, -0.5) -- (1.2, -0.5) node[right, font=\tiny, black] {Air neuf / soufflé};
    \draw[flux_e] (3.5, -0.5) -- (4.7, -0.5) node[right, font=\tiny, black] {Air extrait / rejeté};
    \draw[FedGreen, thick, dashed, ->] (7.0, -0.5) -- (8.2, -0.5)
      node[right, font=\tiny, black] {Free-cooling (bypass)};

  \end{tikzpicture}
  \caption{Schéma de principe CTA double flux avec récupération échangeur rotatif ($\eta \geq 85\,\%$),
    batteries froide et chaude, free-cooling par bypass — Bâtiment tertiaire Lyon, RE2020}
  \label{fig:sr2-schema-cta}
\end{figure}

% ================================================================
\section{Cahier des charges}
% ================================================================

\begin{methode}[Livrables attendus — Note 1 : Dossier technique écrit]
\textbf{À remettre en semaine 6 — Format PDF, 30 pages minimum}

\begin{enumerate}
  \item \textbf{Bilan aéraulique} :
    \begin{itemize}
      \item Calcul des débits d'air neuf et de reprise selon NF EN 16798 et arrêté du 24/03/1982
      \item Dimensionnement du réseau principal de gaines (vitesses, sections, pertes de charge)
      \item Calcul du SFP (Specific Fan Power) en \si{\watt\per\litre\per\second}
      \item Vérification de conformité avec la limite SFP RE2020
    \end{itemize}

  \item \textbf{Sélection et dimensionnement de la CTA} :
    \begin{itemize}
      \item Calcul des puissances thermiques (batterie froide, batterie chaude)
      \item Calcul de la récupération de chaleur (gains en kWh/an)
      \item Sélection d'une CTA constructeur (Airmaster, Carrier, Daikin) avec justification
      \item Calcul des économies de free-cooling (heures/an sur données météo Lyon)
    \end{itemize}

  \item \textbf{Programme GTB} :
    \begin{itemize}
      \item Liste des points d'entrée/sortie (capteurs, actionneurs)
      \item Stratégie de régulation : boucles PID, consignes saisonnières
      \item Scénario de free-cooling : conditions de déclenchement et de rebasculement
      \item Planning horaire type (semaine, week-end, vacances)
      \item Architecture réseau (KNX, BACnet ou Modbus — à justifier)
    \end{itemize}

  \item \textbf{Bilan énergétique RE2020} :
    \begin{itemize}
      \item Calcul du Cep\textsubscript{nr} de référence et du projet [\si{\kilo\watt\hour_{ep}\per\metre\squared\per\year}]
      \item Calcul du Bbio (indicateur besoin bioclimatique)
      \item Calcul de l'Ic\textsubscript{énergie} (impact carbone)
      \item Vérification des 3 seuils RE2020
    \end{itemize}

  \item \textbf{Démarche économique} :
    \begin{itemize}
      \item Coût d'installation estimatif (CTA, gaines, régulation, GTB)
      \item Économies annuelles vs solution sans récupération ni free-cooling
      \item Retour sur investissement
    \end{itemize}
\end{enumerate}
\end{methode}

\begin{methode}[Livrables attendus — Note 2 : Soutenance + démonstration GTB]
\textbf{Durée : 20 min exposé + 10 min questions — Semaine 7}

\begin{itemize}
  \item Présentation des choix techniques face à 2 alternatives (ex : VRV vs CTA centralisée)
  \item \textbf{Démonstration GTB obligatoire :} simuler en direct sur tableur (ou logiciel) :
    \begin{itemize}
      \item Scénario 1 : déclenchement free-cooling ($T_{ext} = 18°C$, $T_{int} = 24°C$)
      \item Scénario 2 : alarme filtre colmaté (déséquilibre pression)
    \end{itemize}
  \item Questions individuelles sur les calculs de bilan aéraulique
\end{itemize}
\end{methode}

% ================================================================
\section{Grille d'évaluation}
% ================================================================

\begin{center}
\begin{tabular}{p{5.5cm}ccp{4.5cm}}
  \toprule
  \textbf{Critère} & \textbf{Barème} & \textbf{Note} & \textbf{Indicateurs} \\
  \midrule
  \multicolumn{4}{l}{\textbf{Note 1 — Dossier technique (20 pts)}} \\
  \midrule
  Bilan aéraulique (débits, $\Delta P$, SFP) & 5 pts & /5 & Conformité NF EN 16798 \\
  Dim. CTA + récupération & 4 pts & /4 & Puissances cohérentes, modèle justifié \\
  Programme GTB & 4 pts & /4 & Points E/S complets, scénario free-cooling \\
  Bilan RE2020 & 4 pts & /4 & Cep\textsubscript{nr} $\leq 50$, 3 indicateurs calculés \\
  Qualité du dossier & 3 pts & /3 & Schémas légendés, rédaction pro \\
  \midrule
  \multicolumn{4}{l}{\textbf{Note 2 — Soutenance + démo GTB (20 pts)}} \\
  \midrule
  Exposé technique & 5 pts & /5 & Clarté, structure, gestion temps \\
  Démonstration GTB & 6 pts & /6 & Scénarios corrects, réaction juste \\
  Questions individuelles & 6 pts & /6 & Maîtrise des calculs propres \\
  Argumentation des choix & 3 pts & /3 & Comparaison alternatives \\
  \bottomrule
\end{tabular}
\end{center}

% ================================================================
\section{Ressources documentaires}
% ================================================================

\begin{itemize}
  \item \textbf{Normes :} NF EN 16798 (aéraulique tertiaire), NF EN 15232 (GTB), NF EN 13053 (CTA)
  \item \textbf{Réglementation :} RE2020 tertiaire, Arrêté du 24/03/1982 (aération), Décret tertiaire 2019
  \item \textbf{Logiciels :} ClimaWin (bilan aéraulique), e-COSTIC (GTB), EnerCalc (RE2020)
  \item \textbf{Catalogues :} Airmaster Duplex, Carrier AHU, Daikin Modular
  \item \textbf{Cours :} Chapitres 03 (aéraulique), 06 (climatisation), 07 (GTB/régulation), 08 (RE2020)
\end{itemize}

\finchapitre


% ================================================================
%  Situation Réelle 3 — Conception complète CVC d'un immeuble résidentiel R+4
%  BTS FED — Option GCF — Semestre 3 / Fin d'année
%  Chapitres mobilisés : TOUS (00 à 11) — Projet transversal
% ================================================================

\chapter{SR-3 — Conception complète CVC d'un immeuble résidentiel R+4}

\begin{tcolorbox}[
  enhanced, colback=FedGreen!8, colframe=FedGreen, arc=3pt,
  fonttitle=\bfseries\sffamily\large, coltitle=white, colbacktitle=FedGreen,
  title={Situation Réelle 3 — Fiche d'identité du projet},
  width=\textwidth
]
\begin{tabular}{@{}ll@{\hspace{2cm}}ll@{}}
  \textbf{Semestre} & 3 (fin d'année) &
  \textbf{Durée} & 8 semaines \\[4pt]
  \textbf{Groupe} & 2 à 3 étudiants &
  \textbf{Chapitres} & 00 à 11 (transversal) \\[4pt]
  \textbf{Note 1} & Dossier complet (coeff.\ 3) &
  \textbf{Note 2} & Soutenance commission (coeff.\ 2) \\[4pt]
  \textbf{Compétences} & \textbf{Toutes} C1 à C9 & & \\
\end{tabular}
\end{tcolorbox}

\begin{attention}[Projet de fin d'année — Niveau professionnel attendu]
  Ce projet constitue le \textbf{projet de synthèse} de la formation BTS FED GCF.
  Le dossier technique doit atteindre un niveau proche d'une note de calcul professionnelle
  (style bureau d'études), avec plan de principe, notes de calcul justifiées, sélections
  de matériels sur catalogues réels et vérification réglementaire complète.
\end{attention}

\bigskip

% ================================================================
\section{Contexte professionnel}
% ================================================================

Vous êtes ingénieur d'études chez \textbf{FluidEnergy Bureau d'Études} (Bordeaux, 33).
Un promoteur, \textbf{Gironde Habitat}, vous confie la conception complète des installations
CVC, VMC et électriques CVC d'un \textbf{immeuble résidentiel R+4} neuf,
situé à \textbf{Bordeaux Bastide} (zone H2a, latitude 44,8°N).

Le programme impose des contraintes fortes sur le choix des équipements et le respect
de la RE2020 dans sa version résidentielle.

\subsection*{Description architecturale de l'immeuble}

\begin{figure}[htbp]
  \centering
  \begin{tikzpicture}[scale=0.9, >=Stealth, thick]

    % ============ COUPE TRANSVERSALE ============
    \node[font=\bfseries\small, FedBlue] at (6, 12.8) {Coupe transversale — Immeuble R+4 Gironde Habitat};

    % === Sous-sol parking ===
    \fill[gray!30] (0, -2.5) rectangle (12, 0);
    \draw[gray!60, thick] (0,-2.5) rectangle (12,0);
    \node[font=\small\bfseries, gray!70] at (6,-1.25) {Sous-sol — Parking 15 places};
    % PAC géo dans SS
    \draw[FedGreen, thick, fill=FedGreen!20, rounded corners=2pt]
      (0.5,-2.2) rectangle (3.0,-0.5);
    \node[FedGreen, font=\tiny\bfseries] at (1.75,-1.0) {PAC géoth.};
    \node[FedGreen, font=\tiny] at (1.75,-1.4) {2×25 kW};
    \node[FedGreen, font=\tiny] at (1.75,-1.7) {cascadées};
    % Sonde géo
    \draw[FedGreen, thick, ->] (1.75,-2.2) -- (1.75,-3.2)
      node[below, font=\tiny] {Sondes \SI{100}{\metre}};
    % Local technique SS
    \draw[FedBlue!60, thick, fill=FedBlue!10, rounded corners=2pt]
      (9.0,-2.2) rectangle (11.5,-0.5);
    \node[FedBlue, font=\tiny\bfseries] at (10.25,-1.0) {Local tech.};
    \node[FedBlue, font=\tiny] at (10.25,-1.4) {Chaufferie};
    \node[FedBlue, font=\tiny] at (10.25,-1.7) {+ Nourrice};

    % === RDC ===
    \fill[FedBlue!8] (0,0) rectangle (12,2.5);
    \draw[FedBlue!60] (0,0) rectangle (12,2.5);
    % Entrée + local vélos + poussettes
    \draw[brown!60, fill=brown!15] (4.5,0) rectangle (7.5,2.5);
    \node[brown!80, font=\tiny] at (6,1.25) {Hall + locaux communs};
    % 2 logements RDC
    \draw[FedBlue!40, fill=white] (0.2,0.2) rectangle (4.3,2.3);
    \node[FedBlue, font=\tiny] at (2.25,1.25) {T3 — \SI{62}{\metre\squared}};
    \draw[FedBlue!40, fill=white] (7.7,0.2) rectangle (11.8,2.3);
    \node[FedBlue, font=\tiny] at (9.75,1.25) {T3 — \SI{62}{\metre\squared}};
    \node[right, font=\tiny\bfseries, FedBlue] at (12.1,1.25) {RDC};

    % === R+1 à R+4 ===
    \foreach \niveau/\label in {2.5/R+1, 5.0/R+2, 7.5/R+3, 10.0/R+4}{
      \fill[FedBlue!5] (0,\niveau) rectangle (12,\niveau+2.3);
      \draw[FedBlue!40] (0,\niveau) rectangle (12,\niveau+2.3);
      % 4 logements par niveau
      \draw[FedBlue!30, fill=white] (0.2,\niveau+0.2) rectangle (2.8,\niveau+2.1);
      \draw[FedBlue!30, fill=white] (3.0,\niveau+0.2) rectangle (5.6,\niveau+2.1);
      \draw[FedBlue!30, fill=white] (6.4,\niveau+0.2) rectangle (9.0,\niveau+2.1);
      \draw[FedBlue!30, fill=white] (9.2,\niveau+0.2) rectangle (11.8,\niveau+2.1);
      \node[FedBlue, font=\tiny] at (1.5,\niveau+1.15) {T2};
      \node[FedBlue, font=\tiny] at (4.3,\niveau+1.15) {T3};
      \node[FedBlue, font=\tiny] at (7.7,\niveau+1.15) {T3};
      \node[FedBlue, font=\tiny] at (10.5,\niveau+1.15) {T2};
      \node[right, font=\tiny\bfseries, FedBlue] at (12.1,\niveau+1.15) {\label};
    }

    % === Toiture ===
    \fill[gray!20] (0,12.3) rectangle (12,12.8);
    \draw[gray, thick] (0,12.3) -- (12,12.3);
    % CTA toiture
    \draw[FedOrange, thick, fill=FedOrange!15, rounded corners=2pt]
      (4.5,12.3) rectangle (7.5,12.8);
    \node[FedOrange, font=\tiny\bfseries] at (6,12.55) {CTA VMC collective};

    % === Colonnes montantes (tirets) ===
    \draw[FedOrange, thick, dashed] (3.2, -0.5) -- (3.2, 12.3);
    \node[FedOrange, font=\tiny, rotate=90] at (3.35, 6) {Colonne chauffage};
    \draw[FedBlue, thick, dashed] (8.8, -0.5) -- (8.8, 12.3);
    \node[FedBlue, font=\tiny, rotate=90] at (8.65, 6) {Colonne ECS};

    % === Légende ===
    \draw[FedGreen, thick, fill=FedGreen!20] (0.2,-4.2) rectangle (1.2,-3.7);
    \node[right, font=\tiny] at (1.3,-3.95) {PAC géothermique};
    \draw[FedOrange, thick, dashed] (4.0,-3.95) -- (5.0,-3.95);
    \node[right, font=\tiny] at (5.1,-3.95) {Colonne chauffage};
    \draw[FedBlue, thick, dashed] (8.0,-3.95) -- (9.0,-3.95);
    \node[right, font=\tiny] at (9.1,-3.95) {Colonne ECS};

  \end{tikzpicture}
  \caption{Coupe transversale de l'immeuble R+4 Gironde Habitat — Bordeaux (H2a) —
    20 logements (2×T2 + 2×T3 par niveau R+1 à R+4, 2×T3 au RDC) —
    PAC géothermique en sous-sol + colonnes montantes + VMC hygroB}
  \label{fig:sr3-coupe-immeuble}
\end{figure}

\subsection*{Données du programme immobilier}

\begin{center}
\begin{tabular}{lrl}
  \toprule
  \textbf{Caractéristique} & \textbf{Valeur} & \textbf{Remarque} \\
  \midrule
  Nombre de logements & 20 & 4 T2 + 16 T3 (voir coupe) \\
  Surface totale SHON & \SI{1380}{\metre\squared} & RE2020 résidentiel \\
  Zones climatiques & H2a & $T_{ext,base} = \SI{-5}{\celsius}$ \\
  Hauteur sous plafond & \SI{2.5}{\metre} & Logements / \SI{3.0}{\metre} SS \\
  Sous-sol parking & 15 places & Ventilation mécanique imposée \\
  \midrule
  \textbf{Contraintes maître d'ouvrage} & & \\
  \midrule
  Production chaleur & PAC géothermique sondes & 2 unités cascadées \\
  Chauffage & Plancher chauffant basse T° & Chaque logement \\
  ECS & Chauffe-eau thermodynamique & 1 par logement \\
  VMC & Hygroréglable de type B & NF EN 13141-8 \\
  Tableau électrique CVC & 1 armoire dédiée par cage & Bilan de puissance obligatoire \\
  \bottomrule
\end{tabular}
\end{center}

% ================================================================
\section{Schéma général de l'installation}
% ================================================================

\begin{figure}[htbp]
  \centering
  \begin{tikzpicture}[
    >=Stealth, thick,
    bloc/.style={draw, rounded corners=3pt, minimum width=2.5cm, minimum height=1.0cm,
                 align=center, font=\scriptsize},
    dep/.style={FedOrange, thick, ->},
    ret/.style={FedBlue, thick, ->},
    geo/.style={FedGreen, thick, ->},
  ]

    % === SOURCES (gauche) ===
    \node[bloc, fill=FedGreen!20, draw=FedGreen, minimum height=1.4cm] (pac) at (0,4)
      {PAC géo\\2×25 kW\\COP = 4,2};
    \node[bloc, fill=brown!15, draw=brown!60] (sondes) at (0,1.5)
      {Sondes\\géothermiques\\2×100 m};
    \draw[geo, ->] (sondes.north) -- node[left, font=\tiny] {$T_{sol} = 13°C$} (pac.south);

    % === CIRCUIT PRIMAIRE ===
    \node[bloc, fill=FedBlue!10, draw=FedBlue] (tampon) at (4,4)
      {Ballon\\tampon\\500 L};
    \draw[dep] (pac.east) -- node[above, font=\tiny] {45°C} (tampon.west);
    \draw[ret] (tampon.west|-tampon.south) -- node[below, font=\tiny] {35°C}
      (pac.east|-pac.south);

    % === NOURRICE / COLLECTEUR ===
    \node[bloc, fill=FedOrange!15, draw=FedOrange, minimum width=3.0cm] (nourrice) at (8,4)
      {Nourrice\\départ/retour};
    \draw[dep] (tampon.east) -- (nourrice.west);
    \draw[ret, bend right=20] (nourrice.south) to (tampon.east|-tampon.south);

    % === CIRCUITS SECONDAIRES ===
    % PCH logements
    \node[bloc, fill=gray!10] (pch) at (12, 6)
      {Colonnes PCH\\20 logements\\$\phi = 25$°C / 35°C};
    % ECS
    \node[bloc, fill=cyan!10, draw=cyan!60] (ecs) at (12, 4)
      {CET\\20 logements\\200 L / logement};
    % Ventilation parking
    \node[bloc, fill=yellow!15, draw=yellow!60] (parking) at (12, 2)
      {VMC parking\\SS — \SI{3000}{\metre\cubed\per\hour}};

    \draw[dep, ->] (nourrice.north east) |- (pch.west);
    \draw[dep, ->] (nourrice.east) -- (ecs.west);
    \draw[dep, ->] (nourrice.south east) |- (parking.west);

    % Retours
    \draw[ret, ->] (pch.south) -- ++(0,-0.5) -| (nourrice.north);
    \draw[ret, ->] (ecs.south) -- ++(0,-0.3) -| (nourrice.north|-nourrice.south);
    \draw[ret, ->] (parking.north) -- ++(0,0.3) -| (nourrice.south);

    % === VASE D'EXPANSION ===
    \node[draw, circle, minimum size=0.7cm, fill=yellow!20, font=\tiny] (vase) at (4,6.5)
      {VE};
    \draw[gray, dashed] (vase) -- (tampon.north);
    \node[font=\tiny, gray] at (4,7.0) {25 L};

    % === COMPTAGE ===
    \node[draw, diamond, fill=purple!10, font=\tiny, minimum size=0.8cm] (compt) at (8,6.5)
      {Comptage\\calorimétrique};
    \draw[purple, dashed] (compt) -- (nourrice.north);

    % === RÉGULATION ===
    \node[bloc, fill=purple!10, draw=purple!50, minimum width=2.0cm] (reg) at (0,7.5)
      {Régulateur\\GTB\\BACnet IP};
    \draw[purple, dashed, ->] (reg.south) -- (pac.north);
    \draw[purple, dashed, ->] (reg.east) -| (nourrice.north|-compt) -- (compt.west);

    % === LÉGENDE ===
    \draw[dep] (0,-1.0) -- (1.2,-1.0) node[right, font=\tiny, black] {Départ chaud};
    \draw[ret] (3.5,-1.0) -- (4.7,-1.0) node[right, font=\tiny, black] {Retour froid};
    \draw[geo] (7.0,-1.0) -- (8.2,-1.0) node[right, font=\tiny, black] {Circuit géothermique};
    \draw[purple, dashed] (10.5,-1.0) -- (11.7,-1.0) node[right, font=\tiny, black] {GTB / Régulation};

  \end{tikzpicture}
  \caption{Schéma général de l'installation CVC — Immeuble R+4 Bordeaux —
    PAC géothermique 2×25 kW + ballon tampon 500 L + nourrice + colonnes PCH + CET + VMC parking}
  \label{fig:sr3-schema-installation}
\end{figure}

% ================================================================
\section{Cahier des charges}
% ================================================================

\begin{methode}[Livrables attendus — Note 1 : Dossier de conception complet]
\textbf{À remettre en semaine 8 — Minimum 50 pages — Niveau note de calcul professionnelle}

\begin{enumerate}
  \item \textbf{Analyse dimensionnelle et cohérence des calculs} (ch. 00) :
    \begin{itemize}
      \item Vérification dimensionnelle des 3 formules principales utilisées dans le dossier
      \item Tableau de synthèse des grandeurs, unités SI et ordres de grandeur attendus
    \end{itemize}

  \item \textbf{Bilan thermique complet} (ch. 01) :
    \begin{itemize}
      \item Calcul par type de logement (T2, T3) selon NF EN 12831
      \item $U_{bat}$ et comparaison seuil RE2020 résidentiel
      \item Déperditions totales bâtiment [\si{\kilo\watt}]
    \end{itemize}

  \item \textbf{PAC géothermique} (ch. 04) :
    \begin{itemize}
      \item Dimensionnement puissance, sélection modèle, calcul COP et SPF annuel
      \item Dimensionnement des sondes géothermiques (longueur, résistance thermique du sol)
      \item Vérification compatibilité avec le plancher chauffant (T° départ $\leq 45°C$)
    \end{itemize}

  \item \textbf{Réseau hydraulique} (ch. 02) :
    \begin{itemize}
      \item Dimensionnement des colonnes montantes et antennes de logement
      \item Équilibrage hydraulique par vannes d'équilibrage (méthode index + $k_{v}$)
      \item Vase d'expansion DTU 65.11 + sécurités (soupape, pressostat)
      \item Sélection des circulateurs (ErP 2015 + adaptation vitesse)
    \end{itemize}

  \item \textbf{Plancher chauffant} (ch. 05) :
    \begin{itemize}
      \item Calcul de la densité de flux $\phi$ [\si{\watt\per\metre\squared}] et du pas de pose
      \item Dimensionnement des boucles (longueur, $\Delta P$) par logement type
      \item Vérification de la température de surface $\leq 29°C$ (NF EN 1264)
    \end{itemize}

  \item \textbf{VMC hygroréglable de type B} (ch. 03) :
    \begin{itemize}
      \item Débits hygroréglables par local (cuisine, salle de bain, WC)
      \item Dimensionnement réseau collectif + gaine verticale en sous-pression
      \item SFP réseau [\si{\watt\per\litre\per\second}] et classe énergétique
    \end{itemize}

  \item \textbf{Étude acoustique CVC} (ch. 09) :
    \begin{itemize}
      \item Niveau de puissance acoustique $L_w$ des équipements (PAC, VMC, circulateurs)
      \item Vérification du niveau NC dans les logements (limite NRA : $\leq 30 dB(A)$ nuit)
      \item Mesures constructives si dépassement (silencieux, plots antivibratoires)
    \end{itemize}

  \item \textbf{Électrotechnique CVC} (ch. 10) :
    \begin{itemize}
      \item Bilan de puissance installée et absorbée [\si{\kilo\watt}]
      \item Schéma unifilaire de l'armoire CVC (disjoncteurs, contacteurs, variateurs)
      \item Calcul de $\cos\varphi$ global et compensation éventuelle
    \end{itemize}

  \item \textbf{Bilan RE2020 complet} (ch. 08) :
    \begin{itemize}
      \item Calcul Bbio, Cep\textsubscript{nr}, Ic\textsubscript{énergie}
      \item Attestation de prise en compte RE2020 (simulation Pleiades ou justification)
      \item Vérification des 3 seuils réglementaires
    \end{itemize}

  \item \textbf{Budget et planning} :
    \begin{itemize}
      \item Devis estimatif par lot (génie thermique, VMC, électricité CVC)
      \item Planning de chantier Gantt sur 8 semaines (gros œuvre terminé)
    \end{itemize}
\end{enumerate}
\end{methode}

\begin{methode}[Livrables attendus — Note 2 : Soutenance devant commission]
\textbf{Durée : 30 min exposé + 15 min questions — Commission mixte : enseignant + professionnel invité}

\begin{itemize}
  \item Plan de principe A3 imprimé à apporter (schéma général de l'installation)
  \item Présentation des choix technologiques et comparaison avec 2 alternatives justifiées :
    \begin{itemize}
      \item Alternative A : PAC géo vs PAC air/eau (COP, investissement, faisabilité)
      \item Alternative B : PCH vs radiateurs basse T° (inertie, confort, coût)
    \end{itemize}
  \item Questions de jury sur \textbf{n'importe quel calcul du dossier} — préparation individuelle obligatoire
  \item Mise en situation : le jury annonce une contrainte de chantier (ex : sondes géo impossibles) — adapter le projet en 5 minutes
\end{itemize}
\end{methode}

% ================================================================
\section{Grille d'évaluation}
% ================================================================

\begin{center}
\begin{tabular}{p{6cm}ccp{4cm}}
  \toprule
  \textbf{Critère} & \textbf{Barème} & \textbf{Note} & \textbf{Indicateurs} \\
  \midrule
  \multicolumn{4}{l}{\textbf{Note 1 — Dossier de conception (20 pts)}} \\
  \midrule
  Bilan thermique + PAC géo & 4 pts & /4 & Puissances cohérentes, SPF calculé \\
  Réseau hydraulique + PCH & 4 pts & /4 & Équilibrage, $\Delta P$, vase \\
  VMC hygroB + acoustique & 3 pts & /3 & SFP, NC logements vérifiés \\
  Électrotechnique CVC & 2 pts & /2 & Armoire, $\cos\varphi$ \\
  RE2020 complet & 3 pts & /3 & 3 indicateurs, seuils respectés \\
  Budget + planning & 2 pts & /2 & Cohérence €, Gantt structuré \\
  Qualité professionnelle & 2 pts & /2 & Mise en page, schémas, vocabulaire \\
  \midrule
  \multicolumn{4}{l}{\textbf{Note 2 — Soutenance commission (20 pts)}} \\
  \midrule
  Exposé (clarté, structure) & 5 pts & /5 & Plan logique, gestion du temps \\
  Maîtrise des calculs (jury) & 7 pts & /7 & Réponses exactes aux questions \\
  Argumentation des choix & 5 pts & /5 & Alternatives comparées, pertinence \\
  Réactivité situation imposée & 3 pts & /3 & Adaptation rapide et cohérente \\
  \bottomrule
\end{tabular}
\end{center}

% ================================================================
\section{Ressources documentaires}
% ================================================================

\begin{itemize}
  \item \textbf{Normes :} NF EN 12831, NF EN 14511 (PAC), NF EN 1264 (PCH), NF EN 13141-8 (VMC hygro B)
  \item \textbf{DTU :} 65.11 (sécurités hydrauliques), 68.1 (VMC), 60.1 (plomberie)
  \item \textbf{Réglementation :} RE2020 résidentiel (arrêté 04/08/2021), Décret F-Gas 2024
  \item \textbf{Logiciels conseillés :} Pleiades+Comfie (RE2020), ClimaWin, Trace 700
  \item \textbf{Catalogues :} Kensa (PAC géo), Grundfos (circulateurs), Zehnder (VMC), Giacomini (PCH)
  \item \textbf{Données météo Bordeaux :} Fichier DRY METEONORM — station 33522 — 2 900 DJU
  \item \textbf{Cours :} Tous les chapitres 00 à 11 du référentiel GCF
\end{itemize}

\finchapitre


\end{document}
